
\subsection{四条 Cauchy--Schwarz 不等式，与冠军的「四重取等」}\label{sec:fourcs}

\begin{proposition}[行方差不等式组，已证]\label{prop:fourcs}
记号同第~\ref{sec:rowid} 小节。则下列四条\textbf{无条件}成立（不需要自洽性）：
\begin{align}
\text{\textbf{[f]}}\quad &(\rho-g_i)^2\,\pi_i\ \le\ a_iv_i
&&\text{（每一行，正部行和的 CS）}\label{eq:f}\\
\text{\textbf{[g]}}\quad &\beta^2\ \le\ S_OW_O
&&\text{（负部的整体 CS）}\label{eq:g}\\
\text{\textbf{[h1]}}\quad &\beta\ \le\ \sum_i\sqrt{\pi_ir_iw_i}
&&\text{（负部的逐行 CS）}\label{eq:h1}\\
\text{\textbf{[h2]}}\quad &\alpha(1-\alpha)+\beta^2\ \le\ R
&&\text{（命题~\ref{prop:perronrig}(c) $+$ 逐行 CS）}\label{eq:h2}
\end{align}
\end{proposition}

\begin{proof}
\eqref{eq:f}：由引理~\ref{lem:rowid}(ii) 与 Cauchy--Schwarz，
$s_i^2=\Bigl(\sum_{(i,j)\in A}(q_{ij})_+x_j\Bigr)^2\le
\Bigl(\sum_{(i,j)\in A}\bigl((q_{ij})_+\bigr)^2\Bigr)\Bigl(\sum_{(i,j)\in A}x_j^2\Bigr)=a_iv_i$，
而 $s_i=(\rho-g_i)x_i$。

\eqref{eq:g}：$\beta=\sum_{(i,j)\in O}(-q_{ij})x_ix_j\le
\sqrt{\bigl(\sum_Oq_{ij}^2\bigr)\bigl(\sum_Ox_i^2x_j^2\bigr)}$，
其中 $\sum_Oq_{ij}^2=\sum_ir_i=S_O$（每个非对角元被有序对数一次），$\sum_O\pi_i\pi_j=\sum_i\pi_iw_i=W_O$。

\eqref{eq:h1}：$d_i=(Q_-x)_i=\sum_{(i,j)\in O}(-q_{ij})x_j\le\sqrt{r_iw_i}$，
故 $\beta=\langle x,d\rangle=\sum_ix_id_i\le\sum_i\sqrt{\pi_ir_iw_i}$。

\eqref{eq:h2}：$\delta=\|d\|^2=\sum_id_i^2\le\sum_ir_iw_i=R$，代入命题~\ref{prop:perronrig}(c) 的
$\delta=\beta^2+\alpha(1-\alpha)$。$\square$
\end{proof}

\begin{remark}[冠军是这四条的\textbf{四重取等}点]\label{rmk:quadtight}
对冠军 $Q^*\in\mathcal P_3(6)$（$K_1\vee C_5$ 核心、权 $w^*=(a^*,b^*\times5)$、定理~\ref{thm:champ}），
数值到 $10^{-16}$：
\begin{center}\small
\begin{tabular}{lrr}
\toprule
量 & 值 & 说明\\
\midrule
$g_{\rm hub},g_{\rm cone}$ & $0.331767104061,\ 0.533646579188$ & 和为 $3$\\
$\pi_{\rm hub},\pi_{\rm cone}$ & $0.197271565265,\ 0.160545686947$ & 和为 $1$\\
$a_{\rm hub},a_{\rm cone}$ & $0.221697692724,\ 0.176698138437$ & \\
$r_{\rm hub},r_{\rm cone}$ & $0,\ 0.072169769272$ & 枢纽行无钝角\\
$v_{\rm hub},v_{\rm cone}$ & $0.802728434735,\ 0.518362939159$ & \\
$w_{\rm hub},w_{\rm cone}$ & $0,\ 0.321091373894$ & \\
\midrule
\eqref{eq:f} 枢纽行 & $0.177963041865=0.177963041865$ & \textbf{取等}（残差 $2.8\times10^{-17}$）\\
\eqref{eq:f} 锥顶点行 & $0.089807129438\le0.091593766384$ & 缺口 $1.787\times10^{-3}$\\
\eqref{eq:g} & $\beta^2=0.093008492800568=S_OW_O$ & \textbf{取等}（残差 $4.2\times10^{-17}$）\\
\eqref{eq:h1} & $\beta=0.304972937817=\sum\sqrt{\pi_ir_iw_i}$ & \textbf{取等}\\
\eqref{eq:h2} & $\alpha(1-\alpha)+\beta^2=0.115865451846=R$ & \textbf{取等}\\
\bottomrule
\end{tabular}
\end{center}
取等的原因很清楚：$Z_5$-对称使枢纽行的 $5$ 个正元全相等、$x$ 在锥顶点上也全相等，
于是 \eqref{eq:f} 的 CS 自动取等；钝角图是五角星、$10$ 个钝角元全相等，
于是 \eqref{eq:g}\eqref{eq:h1}\eqref{eq:h2} 也全部取等。
\textbf{唯一不取等的是 $5$ 个锥顶点行的 \eqref{eq:f}}（缺口 $1.787\times10^{-3}$），
原因同样清楚：锥顶点行有 $1$ 个枢纽正元 $0.210570$ 与 $2$ 个近邻正元 $0.257253$，两者不等，
故 CS 的「成比例」条件失效。\textbf{这 $1.787\times10^{-3}$ 就是全部主线的来源}
（第~\ref{sec:diag} 小节把这句话精确化）。
\end{remark}

\subsection{行方差松弛 REL 与天花板 $L(N)$：Mohar 与 \STW\ 常数的统一重导}\label{sec:relax}

把命题~\ref{prop:fourcs} 与引理~\ref{lem:rowid}(i) 里的 $a_i=g_i(1-g_i)-r_i$ 合起来，
就得到一个\textbf{只含 $(g,\pi,r,w)$ 的有限维松弛}：$Q$ 本身、符号图案、秩三条件统统被丢掉，
只留下「行能量如何瓜分」与「钝角质量如何分布」。

\begin{definition}[行方差松弛 REL]\label{def:rel}
对给定的 $N$，REL$(N)$ 是如下最大化问题的最优值：在
$g\in[0,1]^N$、$\sum g_i=3$；$\pi\ge0$、$\sum\pi_i=1$；$0\le r_i\le g_i(1-g_i)$；$0\le w_i\le1-\pi_i$
上，最大化 $\rho=\alpha+\beta$，约束为 $0\le\alpha\le1$、$\beta\ge0$，以及
\begin{equation}\label{eq:rel}
(\rho-g_i)^2\pi_i\le\bigl(g_i(1-g_i)-r_i\bigr)(1-\pi_i-w_i)\ \ \forall i,\qquad
\beta^2\le S_OW_O,\qquad \beta\le\sum_i\sqrt{\pi_ir_iw_i},\qquad
\alpha(1-\alpha)+\beta^2\le R .
\end{equation}
\end{definition}

由命题~\ref{prop:fourcs} 与引理~\ref{lem:rowid}(i)，任何真实构型都给出 REL 的一个可行点，故
\begin{equation}\label{eq:relvalid}
\boxed{\ \kappa_3(N)\ \le\ \mathrm{REL}(N)\ }\qquad(\text{无条件}).
\end{equation}
内层（给定 $(g,\pi,r,w)$ 后对 $(\alpha,\beta)$ 的）最大化是\textbf{一维且闭式}的：
$\beta$ 的上限是 $\beta_{\max}=\min\{\sqrt{S_OW_O},\ \sum_i\sqrt{\pi_ir_iw_i},\ \sqrt R\}$，
且 $\rho(\beta)=\beta+\min\{1,\ \frac12(1+\sqrt{(1-4(R-\beta^2))_+}),\ \min_i(g_i+\sqrt{a_iv_i/\pi_i})-\beta\}$
在 $[0,\beta_{\max}]$ 上单调增，故内层最优值就是
\begin{equation}\label{eq:inner}
\lambda^*(g,\pi,r,w)=\min\Bigl\{\,1+\beta_{\max},\ \ \min_i\Bigl(g_i+\sqrt{\tfrac{a_i(1-\pi_i-w_i)}{\pi_i}}\Bigr)\Bigr\}
\quad(\text{当 }\beta_{\max}\text{ 使 }\alpha=1\text{ 可行时}).
\end{equation}
\textbf{冠军处 $\lambda^*=\rho_0$ 精确成立}（数值到 $10^{-16}$；因为 \eqref{eq:f} 的枢纽行取等、
且 $\beta_{\max}=\beta$）。也就是说 REL 在冠军身上\textbf{一点都不松}——松的是别处。

\begin{theorem}[均匀层的天花板 $L(N)$，已证]\label{thm:LN}
设 $Q\in\mathcal P_3(N)$、$Q_+x=\rho x$，且\textbf{双均匀}：$g_i\equiv\frac3N$、$\pi_i\equiv\frac1N$。
记
\[
\kappa_N:=\sqrt{\frac{3(N-1)}{N-3}}\qquad(N\ge4),\qquad
L(N):=\frac{2+\kappa_N}{1+\kappa_N}.
\]
则 $\rho\le L(N)$，且等号成立的必要条件是钝角度
$n^*:=N\omega^*=\dfrac{N\kappa_N}{3(1+\kappa_N)}$（其中 $\omega^*=\dfrac{\kappa_N}{3(1+\kappa_N)}$）。
具体地
\[
L(3)=1,\quad L(4)=\frac54,\quad L(5)=\frac{4+\sqrt6}5=1.289897949,\quad
L(6)=\frac{3+\sqrt5}4=\Phi=1.309016994,
\]
\[
L(7)=\frac{5+3\sqrt2}7=1.320377241,\quad
L(8)=\frac{11+\sqrt{105}}{16}=1.327934423,\quad
L(12)=\frac{5+\sqrt{33}}8=1.343070331,
\]
$L(N)$ 关于 $N$ 严格递增，且 $\displaystyle\lim_{N\to\infty}L(N)=\frac{2+\sqrt3}{1+\sqrt3}=\frac{1+\sqrt3}2=1.366025404$。
\end{theorem}

\begin{remark}[这条定理为什么值得单独立一条]\label{rmk:LNread}
(i) \textbf{Mohar 界与 \STW\ 界是同一个闭式的两端。}
$\lim_{N\to\infty}L(N)=\frac{1+\sqrt3}2$ 正是 Mohar 的常数；$L(6)=\frac{3+\sqrt5}4=\Phi$ 正是 \STW\ 的常数。
两者在我们的框架里\textbf{不是两条独立的不等式}，而是「行方差机制」在 $N=\infty$ 与 $N=6$ 处的取值。
这解释了第~\ref{sec:stwanat} 节留下的疑问：为什么 \STW\ 的 $\Phi$ 偏偏是 $\frac{3+\sqrt5}4$
——因为 $\kappa_6=\sqrt5$，而 $\sqrt5$ 来自 $3(N-1)/(N-3)$ 在 $N=6$ 处的值，即「行能量 $3$」与「$N-3=\dim\ker Q$」。

(ii) $L(4)=\frac54=\kappa_3(4)$（报告已知的精确值），而且\textbf{取等}：
$N=4$ 时 $n^*=\frac{4\cdot3}{3\cdot4}=1$ 是整数，极值构型正是定理~\ref{thm:bipglobal} 的
$K_{2,2}$ 见证 $Q=\frac34I_4+\frac14A(K_{2,2})-\frac14M=I_4-vv^\top$，$v=(\frac12,\frac12,-\frac12,-\frac12)$（谱 $\{1,1,1,0\}$，见注~\ref{rmk:n4exact} 的精确计算）。
所以这个机制\textbf{不是永远松的}：在 $N=4$ 它精确。

(iii) 均匀层只是 REL 的一个子类。我们对 REL 的\emph{自由}（$4N$ 维）上确界做了全局搜索
（差分进化，$170$--$350$ 代，Sobol 初始化，多起点），结果与 $L(N)$ 逐位吻合：
\begin{center}\small
\begin{tabular}{c|cccccc}
\toprule
$N$ & 5 & 6 & 7 & 8 & 10 & 12\\
\midrule
REL 自由搜索 & $1.2898973$ & $1.3090164$ & $1.3203769$ & $1.3279311$ & $1.3373855$ & $1.3430693$\\
$L(N)$ 闭式 & $1.2898979$ & $1.3090170$ & $1.3203772$ & $1.3279344$ & $1.3373864$ & $1.3430703$\\
差 & $-6.5\!\times\!10^{-7}$ & $-5.6\!\times\!10^{-7}$ & $-3.1\!\times\!10^{-7}$ & $-3.4\!\times\!10^{-6}$ & $-9.0\!\times\!10^{-7}$ & $-1.0\!\times\!10^{-6}$\\
\bottomrule
\end{tabular}
\end{center}
（\textbf{数值}标签：这一行的吻合度强烈提示 $\mathrm{REL}(N)=L(N)$，但我们\textbf{只对均匀子类给出了证明}。
若 $\mathrm{REL}(N)=L(N)$ 对一切 $N$ 成立，则 $\kappa_3(N)\le L(N)$ 成为无条件定理，
它在 $N\le6$ 优于或等于 \STW，在 $N\ge7$ 弱于 \STW。）
\end{remark}

\begin{proof}[定理~\ref{thm:LN} 的证明]
双均匀时令 $\varrho:=r_i$（各行的负部能量，由对称性可设相同）、$\omega:=w_i$、$c:=N\varrho$。则
$a_i=\frac3N(1-\frac3N)-\varrho$，$v_i=1-\frac1N-\omega$，
\[
S_O=c,\qquad W_O=\sum_i\pi_iw_i=\omega,\qquad R=c\omega,\qquad
\sum_i\sqrt{\pi_ir_iw_i}=N\sqrt{\tfrac1N\varrho\omega}=\sqrt{c\omega},
\]
即 $\beta_{\max}=\min\{\sqrt{c\omega},\sqrt{c\omega},\sqrt{c\omega}\}=\sqrt{c\omega}$（三个上限\emph{退化成一个}）。
于是 \eqref{eq:inner} 给出 $\rho\le\min\{A,B\}$，其中
\[
A:=\frac3N+\sqrt{\Bigl(3-\frac9N-c\Bigr)\Bigl(1-\frac1N-\omega\Bigr)},\qquad
B:=1+\sqrt{c\omega}.
\]
$A$ 关于 $c$ 严格减、$B$ 关于 $c$ 严格增（$\omega$ 固定时同理反向），故 $\min\{A,B\}$ 的上确界在
$A=B$ 处取得（边界 $c=0$ 或 $\omega=0$ 给 $B=1$，边界 $c=3-\frac9N$ 给 $A=\frac3N$，都更小）。
设 $A=B=L$，记 $K=3-\frac9N=\frac{3(N-3)}N$、$M=1-\frac1N=\frac{N-1}N$，则
\[
(K-c)(M-\omega)=\Bigl(L-\frac3N\Bigr)^2,\qquad c\omega=(L-1)^2 .
\]
把第二式代入第一式并展开，$L^2$ 相消，得
\[
\frac{4(N-3)}N=\frac{3(N-3)}N\,\omega+\frac{N-1}N\,c+\frac{2(N-3)}N\,L,
\]
其中用到 $KM+1-\frac9{N^2}=\frac{3(N-3)(N-1)+N^2-9}{N^2}=\frac{4N^2-12N}{N^2}=\frac{4(N-3)}N$。两边乘 $\frac{N}{2(N-3)}$：
\begin{equation}\label{eq:Lstar}
L=2-\frac32\,\omega-\frac{N-1}{2(N-3)}\,c .
\end{equation}
对 \eqref{eq:Lstar} 的后两项用 AM--GM，并用 $c\omega=(L-1)^2$：
\[
\frac32\omega+\frac{N-1}{2(N-3)}c\ \ge\ 2\sqrt{\frac{3(N-1)}{4(N-3)}\,c\omega}
\ =\ \kappa_N(L-1)\qquad(L\ge1),
\]
代入 \eqref{eq:Lstar} 得 $2-L\ge\kappa_N(L-1)$，即 $L\le\frac{2+\kappa_N}{1+\kappa_N}=L(N)$。
AM--GM 取等当且仅当 $\frac32\omega=\frac{N-1}{2(N-3)}c$，解出
$\omega^*=\frac{\kappa_N}{3(1+\kappa_N)}$、$c^*=\frac{(N-3)\kappa_N}{(N-1)(1+\kappa_N)}$
（直接代入验证 $c^*\omega^*=(L(N)-1)^2$ 成立，故该点确实可行）。
单调性与 $N\to\infty$ 极限由 $\kappa_N$ 关于 $N$ 严格减、$\kappa_N\to\sqrt3$ 得到；
$L(3)$ 单独说明：$N=3$ 时 $g_i\equiv1$，故 $a_i+r_i=g_i(1-g_i)=0$，$Q=I_3$，$\rho=1$。$\square$
\end{proof}
